\section{Conclusiones}
    De las imágenes consideradas para el proyecto, me permitieron 
    darme cuenta de que en la vida real está llena de muchas y diferentes 
    texturas que para estos filtros pueden verse como si fueran 
    bordes, por ello se vuelve vital aplicar una técnica de suavizado 
    para reducir el efecto o implicaciones de considerar estos 
    cambios menores en las tonalidades. Por ello, después del suavizado, 
    los algoritmos como Prewitt, Sobel y Cany generan mejores bordes, 
    debido a que la diferencia entre colores cercanos se redujo o 
    es menor. El problema de los filtros de Prewitt y Sobel es que 
    no son sensibles a las curvas, es decir, se les dificulta detectar 
    elementos no rectos, pero esta es la principal ventaja que tienen 
    los filtros de Laplace y Laplaciano de Gaussiana, son muy sensibles 
    a los cambios de tonalidades (pues comparan vecinos adyacentes de 
    píxeles) por ello pueden detectar curvas muy 
    cerradas con una única pasada, es decir, no se requiere de dos 
    ``versiones'' del mismo filtro sino con una única versión se 
    puede detectar cualquier curva. Por esta misma sensibilidad, es que 
    cualquier cambio en la textura o ruido se verá reflejado inmediatamente  
    como un borde, para sanar esto podríamos introducir un umbral para 
    determinar qué sí es un borde, y esta es la idea de filtro de Canny 
    que emplea de varios filtros direccionales para detectar la mayoría 
    de los bordes en una imagen (sin importar si son rectos o curvos). 
    Por cómo se define, el filtro de Canny es el más complejo de usar 
    y, personalmente, en el proyecto no pude usarlo debidamente porque 
    no pude ajustar los umbrales adecuadamente sin que se descartara 
    un borde o se introducieran no bordes (ruido).\\

    Este miniproyecto se me hizo interesante el como existen métodos 
    no tan sofisticados que pueden ser usados en etapas de preprocesamiento 
    para mejorar la detección, localización y segmentación de objetos 
    en tareas de visión por computadora empleando redes neuronales. 
    Dicho así, no todo se trata de un modelo de IA que lo haga todo, sino 
    también hay que entender y analizar el problema desde otras 
    perspectivas o métodos que puedan integrarse a estos modelos. 
    Y esto es lo que más me llevo de este proyecto, de realizar un paso 
    hacia atrás para comprender un problema desde otra perspectiva 
    empleando técnicas, métodos o modelos más sencillos/intuitivos 
    y explicables (que no sean una caja negra) y que como científicos 
    de datos podamos sacarle toda el conocimiento posible para 
    seguir mejorando la solución final al problema que queremos 
    resolver, evitando dejar huecos teóricos o sin considerar.